\section{Congruence denominators, the good cut, and the original family}
\label{sec:congruence-chain}

\status{LN1--LN10 close the chain from an integral good-cut operator and
explicit local denominator estimates to \texttt{LinkN} on the original
family.  Two useful simplifications are proved: a regular-representation
trace makes the central coefficients integral in the original DVR, and
a filtration by individual eigenspaces bounds a congruence module by its
coordinate congruence ideals.  The argument retains the actual local
lattices, theta ranks, place weights and cohomological sum.  The
per-prime Galois bounds and the compatibility of the actual theta/type
and local data are explicit inputs; their existence is not inferred
from this numerical and algebraic closure.}

Fix one source datum $D$, and write
\[
 \begin{gathered}
 N_G=D.\mathrm{what}_2.\mathrm{restRowsFiniteGroupCard}>0,\\
 Y_i=D.\mathrm{nd}.Z_{\rm good}(i)\subseteq
 Z_i=D.\mathrm{nd}.Z(i)\subseteq X_i.
 \end{gathered}
 \tag{LN.1}
\]
For a localisation $i$, put $\mathcal O=\mathcal O_i$, $K=\operatorname{Frac}
\mathcal O$, $M=M_i$ and $V=K\otimes_{\mathcal O}M$.  The source supplies
a finite set $X=X_i$, a separating $\mathcal O$-subalgebra
$T\subseteq\mathcal O^X$, the actual finite free $\mathcal O$-module
$M$, an integral action of $T$ on $M$, and the
eigenspace-spanning assertion
\[
 V=\sum_{x\in X}V_x,\qquad
 V_x=\{v:f v=f(x)v\text{ for every }f\in T\}.
 \tag{LN.2}
\]
The spanning assertion is a real hypothesis; separation of $T$ alone
does not imply it.  The $\mathcal O_i$ are the original DVRs, finite
over $\mathbb Z_{\ell_i}$ and of characteristic zero through their
specified number-field embeddings.  All lengths below are ordinary
finite module lengths, whose finiteness will be proved before division
or conversion from the source's extended-natural lengths.

\paragraph{LN1: central coefficients in the original coefficient ring.}
\label{claim:LN1}
Let $G$ be a finite group of order $N$, and let
$e\in\mathbb C[G]$ be a central idempotent.  If $\zeta$ is a primitive
$N$th root of unity, then
\[
                   N e(g)\in\mathbb Z[\zeta]\qquad(g\in G).
 \tag{LN.3}
\]
Here $e(g)$ denotes the coefficient of $g$.  To prove this without
character classification, set $y=g^{-1}e$.  Centrality, $e^2=e$ and
$g^N=1$ give $y^{N+1}=y$.  Left multiplication by $y$ on
$\mathbb C[G]$ therefore has eigenvalues in
$\{0\}\cup\{\zeta^a:0\le a<N\}$: its annihilating polynomial is
$t(t^N-1)$.  Its trace is a sum of these eigenvalues and belongs to
$\mathbb Z[\zeta]$.  In the basis $h\in G$, every diagonal entry of
left multiplication by $y$ is $y(1)=e(g)$, so the same trace is
$Ne(g)$.  This proves (LN.3), also for $N=1$.

Suppose the specified embedding $\iota:\mathcal O\hookrightarrow
\mathbb C$ contains a primitive $N$th root.  Lift the coefficients in
(LN.3) through $\iota$ to obtain $z\in\mathcal O[G]$ with
$\iota(z)=Ne$.  The coefficient lifts are unique.  For every $h\in G$,
$hz-zh$ maps to zero, so injectivity gives $hz=zh$.  Thus $z$ belongs
to the center of the original group algebra.  No extension of
$\mathcal O$, change of lattice or unramifiedness assumption is used.

In the actual datum,
\[
 G=\prod_{k=0,1}\prod_{v\in\mathrm{PlG}(k)}
       \mathrm{GL}_r(\mathbb F_{g_k}),\qquad |G|=N_G.
 \tag{LN.4}
\]
The equality is precisely the two source fields
\texttt{card\_PlG} and \texttt{restN\_card\_levelGroup\_eq}.
The chosen \texttt{Pend} supplies a primitive $N_G$th root in every
$\mathcal O_i$.  For every tuple of local labels, the product of the
central local idempotents, embedded in different factors of (LN.4),
is again a central idempotent.  Apply (LN.3) to this product.  This
proves the literal \texttt{CentralElt} coefficient-lifting interface.

The element called $e_0$ in the source is a finite sum over its
\emph{ordered root vectors}.  Lift each good summand as above and add
the lifts, with exactly the original multiplicities.  The resulting
central integral element $z_i$ satisfies
\[
                  \iota(z_i)=N_Ge_0.                 \tag{LN.5}
\]
This proof does not replace the ordered tuples by unordered
polynomials, and does not assert that the sum $e_0$ is an idempotent.
Its specific action on the relevant theta types is the separate
compatibility assertion used next.

\paragraph{LN2: construction of the literal good cut.}
\label{claim:LN2}
We spell out the source inputs used to turn (LN.5) into a cut.  For
each $i$, actual automorphic objects $\Pi$ map to points
$\operatorname{pt}_i(\Pi)\in X_i$ or to no point; every $x\in X_i$
has a representative.  Each such $\Pi$ has a nonzero type subspace.
On every type vector above $x\in Z_i$ the same $e_0$ acts as
\[
 e_0F=F\quad(x\in Y_i),\qquad
 e_0F=0\quad(x\in Z_i\setminus Y_i).                 \tag{LN.6}
\]
This is the exact \texttt{onMt\_e0\_eq\_self\_of\_mem\_zgood\_and\_eq\_zero\_of\_notMem}
interface.  It depends on the actual parameter/type dictionary and
the theta decomposition; LN1 does not prove it.  The original
\texttt{Pend} also gives the integral action $z\mapsto z_{\mathcal O}$,
its commutation with Hecke operators, and \texttt{Pass.EigSpan}: if the
complex action of $z$ is a given scalar on every type space above $x$,
then $(z_{\mathcal O})_K$ is that scalar on $V_x$.

Define, using these actual type spaces,
\[
 S_i=\{x:\exists\Pi,F,\ \operatorname{pt}_i(\Pi)=x,
          \ F\in\operatorname{type}_i(\Pi),\ e_0F\ne0\},
 \qquad \operatorname{keep}(i)\ \Longleftrightarrow\ Y_i\ne\varnothing.
 \tag{LN.7}
\]
Choose a representative and a nonzero type vector for $x\in Y_i$.
Equation (LN.6) proves $Y_i\subseteq S_i$.  Its zero clause proves
$S_i\cap Z_i=Y_i$, and hence $S_i$ is nonempty whenever $i$ is kept.
Let $\alpha_i$ be the action of the actual $z_i$ from (LN.5).
The scalar-compatibility input just stated gives
\[
 [\alpha_i,T_i]=0,\qquad
 (\alpha_i)_K|_{V_x}=N_G\quad(x\in Y_i),\qquad
 (\alpha_i)_K|_{V_x}=0\quad(x\notin S_i).
 \tag{LN.8}
\]
No scalar action on the other selected eigenspaces is required.

There is one further field in \texttt{CutD.Good}: every actual
automorphic representative of a selected point must satisfy
\texttt{Nc} at both generator-prime places.  The precise remaining
type input says that an occurring local idempotent with label $f$
identifies the finite-field inertia-label polynomial with
$f(X)(X-u)^t$, and supplies monodromy zero.  Its degree is the
actual local dimension, at most $n_V$; the generator primes split
completely and satisfy the source's largeness bound, so the
\texttt{Sqf} multiplicity threshold is $2$ there.  For a
good tuple, $f$ is monic, has nonzero constant term, is squarefree,
and $f(u)\ne0$, $f(u^{-1})\ne0$.  Consequently its only possible repeated
label root is $u$ (or $u^{-1}$ at the conjugate place), a unit of
$\mathbb F_{g_k}$ with $(g_k-1)$st power $1$.
The actual complex inertia eigenvalues are read through the
\texttt{Pend} characters $\omega_a(\gamma)$, not by an embedding
of finite fields into $\mathbb C$.  These characters are locally
constant, trivial on wild inertia, separate distinct nonzero label
roots, and are multiplicative in the root $a$.  If two such
characters agree at a tame generator, they agree on its powers
and on wild inertia; density and local constancy give equality on
all inertia, hence equality of their labels.  Evaluation at any
tame generator is therefore injective on the nonzero labels and
preserves their multiplicities.  Also
$a^{g_k-1}=1$ implies $\omega_a(\gamma)^{g_k-1}=1$ by
multiplicativity.  The supplied inertia-characteristic-polynomial
formula now proves exactly the multiplicity-at-least-two clause
of \texttt{Nc}.  On wild inertia every character is $1$, which
proves its \texttt{Tm} clause; with the supplied zero monodromy
this proves \texttt{Nc}.  The argument applies to every tame
generator, as its definition requires.

For completeness, nonzero $e_0F$ supplies the needed occurring
labels.  A nonzero finite sum has a nonzero summand.  If a product of
commuting central local operators is nonzero on $F$, every factor
is nonzero on $F$.  Expanding such a factor as the sum of the
primitive local idempotents having its label produces an occurring
local type with that label.  These are purely finite-dimensional
arguments.  The definition of $S_i$ gives one representative;
the source's \texttt{provisionFortyOne} says that two representatives
of the same selected Hecke point have the same base-change
constituents up to permutation.  Their local parameter multisets
therefore agree.  This proves \texttt{Nc} for \emph{every}
representative, as required, rather than only for the witness in
(LN.7).  Equations (LN.7)--(LN.8) and this last argument construct
all five fields of \texttt{CutD.Good}.

The input boundary is exact: (LN.6), integral/complex action
compatibility, the local parameter dictionary, and the
same-point permutation assertion are retained here.  This is a
complete construction from those fields, not an assertion that
arbitrary finite-group types produce the requisite automorphic data.

\paragraph{LN3: the actual spectral quotient and finite congruence module.}
\label{claim:LN3}
For each $x\in X$, separation gives, for each $y\ne x$, an
$f_{xy}\in T$ with $f_{xy}(x)\ne f_{xy}(y)$.  Then
\[
 h_x=\prod_{y\ne x}(f_{xy}-f_{xy}(y))\in T,
 \quad h_x(x)\ne0,\quad h_x(y)=0\ (y\ne x).          \tag{LN.9}
\]
The empty product is $1$.  The operators $h_x/h_x(x)$ project onto
$V_x$ in (LN.2), so that sum is direct.  Write
$V_A=\bigoplus_{x\in A}V_x$ and $\pi_A:V\to V_A$.
Because $M$ is free, it embeds in $V$; we use that embedding below.

For any $A\subseteq X$, set
\[
 U_A=\pi_A(M)=V_A\cap(M+V_{A^c}),\quad
 L_A=M\cap V_A,\quad C_A(M)=U_A/L_A.                 \tag{LN.10}
\]
The image $U_A$ is a finite free $\mathcal O$-module of rank
$\dim_KV_A$, because it is a finitely generated torsion-free
submodule of a finite-dimensional $K$-space.  Clearing the finitely
many denominators in the projections (LN.9) gives a nonzero
$d\in\mathcal O$ with $dU_A\subseteq L_A$.  Thus $L_A$ has full
rank in $U_A$ and $C_A(M)$ has finite length.  In particular (LN.10)
for $A=Y_i$ is precisely the source's
\texttt{MtUp}/\texttt{MtLow} quotient.

For a selected set $S$, the source defines \texttt{outM} as the
intersection of kernels of the integral operators $f$ which vanish
on $S^c$.  By (LN.9) and the direct sum,
\[
 \operatorname{outM}=M\cap V_{S^c},\qquad
 M'=M/(M\cap V_{S^c}),\qquad
 K\otimes M'\simeq V/V_{S^c}\simeq V_S.              \tag{LN.11}
\]
The quotient $M'$ is torsion-free: if $am\in V_{S^c}$ for a
nonzero $a\in\mathcal O$, then $m\in V_{S^c}$.  Hence $M'$ is
finite free.  Restricting $T$ to $S$ gives the actual \texttt{cutT};
its kernel kills $M'$ and the induced algebra still separates $S$.
The eigenspace at every retained $x$ is the image of $V_x$ under the
quotient, and that map is an isomorphism.  Thus multiplicities of
retained points agree exactly.  These identifications give the
literal source cut module, not a new freely chosen lattice on the
same rational representation.

\paragraph{LN4: a direct filtration bound for coordinate congruence ideals.}
\label{claim:LN4}
For $A\subseteq X$ and $x\in A$, define the nonzero ideal
\[
 I_x(A)=\{f(x):f\in T,\ f|_{A^c}=0\}\subseteq\mathcal O.
 \tag{LN.12}
\]
It is an ideal because $T$ contains constant functions.  It is
nonzero by (LN.9).  The following finite-length estimate is valid
with arbitrary eigenspace multiplicities:
\[
 \operatorname{length}_{\mathcal O} C_A(M)
 \le\sum_{x\in A}\dim_KV_x\,
             \operatorname{length}_{\mathcal O}(\mathcal O/I_x(A)).
 \tag{LN.13}
\]
Order $A=\{x_1,\ldots,x_s\}$ and put
$F_a=\bigoplus_{b\le a}V_{x_b}$ and $U_a=U_A\cap F_a$.
The images of $U_a$ filter $C_A(M)$.  Its $a$th graded quotient is
a quotient of the free module $U_a/U_{a-1}$, whose rank is at most
$\dim_KV_{x_a}$.  Freeness follows because this quotient embeds in
$F_a/F_{a-1}$.  If $f\in T$ vanishes on $A^c$, then
$f\pi_A(m)=fm\in L_A$ for $m\in M$.  On $F_a/F_{a-1}$ it acts by
$f(x_a)$.  More explicitly,
$(f-f(x_a))U_a\subseteq U_{a-1}$, since $f$ preserves $U_A$.
Thus $I_{x_a}(A)$ kills the graded quotient.  A free basis of
$U_a/U_{a-1}$ therefore gives a surjection
$(\mathcal O/I_{x_a}(A))^{\dim V_{x_a}}$ onto it, after adding
zero generators if necessary.  Length additivity along the finite
filtration proves (LN.13).  The case $A=\varnothing$ has both sides
zero.

This argument replaces a longer congruence-module decomposition.
It uses no Galois representation, semisimplicity theorem for an
integral Hecke algebra, or equality of different lattices.  The
ideals (LN.12) are the literal \texttt{ccAt} images: operators
supported on the stated theta subset of the stated $X$.

\paragraph{LN5: cut defect and the original-row denominator inequality.}
\label{claim:LN5}
Suppose $Y\subseteq S$, $S\cap Z=Y$ and $\alpha\in
\operatorname{End}_{\mathcal O}(M)$ obeys (LN.8), with a nonzero
integer $N$ in place of $N_G$.  Let $q:V\to V/V_{S^c}$ be (LN.11).
It induces a map
$C_Y(M)\to C_Y(M')$: indeed $q(U_Y)\subseteq U'_Y$ and
$q(L_Y)\subseteq L'_Y$.  If a class represented by $u\in U_Y$
lies in its kernel, $q(u)=q(m)$ for some $m\in M$; hence
$u-m\in V_{S^c}$ and
\[
              Nu=\alpha u=\alpha m\in M\cap V_Y=L_Y.
 \tag{LN.14}
\]
The kernel is killed by $N$.  Its inverse image in $U_Y$ is a
submodule of a free module of rank $m_Y=\sum_{x\in Y}\dim V_x$,
and is therefore generated by at most $m_Y$ elements.  It follows
that the kernel is a quotient of $(\mathcal O/N\mathcal O)^{m_Y}$.
Length additivity for kernel and image, followed by inclusion of
the image in $C_Y(M')$, gives
\[
 \operatorname{length}C_Y(M)
 \le\operatorname{length}C_Y(M')+
               m_Y\operatorname{length}(\mathcal O/N\mathcal O).
 \tag{LN.15}
\]
Surjectivity onto the cut congruence module is neither needed nor
asserted.

For the normalisation, let $\ell=\ell_i$ and
$e_i=\operatorname{length}_{\mathcal O_i}(\mathcal O_i/\ell\mathcal O_i)$.
This is a positive finite integer.  Indeed $\ell\ne0$, and $\ell$
is not a unit: otherwise the finite nonzero
$\mathbb Z_\ell$-module $\mathcal O_i=\ell\mathcal O_i$ would
vanish by the elementary Nakayama argument (reduce a finite
generating set modulo the maximal ideal, or use its determinant
relation).  The other prime factors of an integer are units through
$\mathbb Z_\ell\to\mathcal O_i$.  Factorising $N$ and adding DVR
lengths gives the exact identity
\[
 \frac{\operatorname{length}(\mathcal O_i/N\mathcal O_i)}{e_i}
                   =v_\ell(N).                         \tag{LN.16}
\]
All involved lengths are finite by LN3 and $N\ne0$, so this is also
the source's \texttt{lenUnit}, including its \texttt{toNat}
convention.  Dividing (LN.15) by $e_i$ proves
\[
 \eta_i\le\eta'_i+
             \sum_{x\in Y_i}\operatorname{mult}_i(x)v_{\ell_i}(N_G).
 \tag{LN.17}
\]
If $Y_i$ is empty, $U_{Y_i}=0$, so $\eta_i=0$; this handles every
discarded localisation.

There is also a useful direct export to the original ideals.  If
$f\in T$ vanishes on $Z^c$, its restriction to $S$ vanishes on
$S\setminus Y$.  Consequently
\[
 I_x(Z)\subseteq I'_x(Y),\qquad
 \operatorname{length}(\mathcal O/I'_x(Y))
       \le\operatorname{length}(\mathcal O/I_x(Z))\quad(x\in Y).
 \tag{LN.18}
\]
Apply LN4 to the cut module and combine (LN.15)--(LN.18).  Define
$v\!c_i(x)=\operatorname{length}(\mathcal O_i/I_x(Z_i))/e_i$.
Then
\[
 \boxed{\quad\eta_i\le
       \sum_{x\in Y_i}\operatorname{mult}_i(x)
           \{v\!c_i(x)+v_{\ell_i}(N_G)\}.\quad}        \tag{LN.19}
\]
This is the actual original-row \texttt{TwistedPlaces.adm1}; the
proof works at every prime where the cut operator is supplied,
including more than its required non-generator range.  The local
denominator $v\!c_i(x)$ remains the original \texttt{vcc}, not the
denominator for an unrelated representation or quotient.

\paragraph{LN6: multiplicities are paid for by the same theta rank.}
\label{claim:LN6}
Fix an actual lift $j$ and a twisted place $w$.  The same complex
space $M_j$ has an internal direct sum decomposition
$M_j=\bigoplus_\sigma E_{j,\sigma}$.  Let $A_i$ be the finite set
of good rows occurring at localisation $i$.  The exact
\texttt{ConclusionZ} identification supplies
\[
 \sum_{x\in Y_i}\operatorname{mult}_i(x)
          \le\sum_{\sigma\in A_i}\dim E_{j,\sigma}
           \quad(i\in\operatorname{locs}(j,w)).          \tag{LN.20}
\]
For a nonempty $Y_i$ this is equality in the source.  For an empty
$Y_i$, the left side is zero and the right side is nonnegative.
\texttt{LocalisationUniqueness} says that the sets $A_i$, for this
fixed $j,w$, are pairwise disjoint.  Summing (LN.20) and using the
internal direct sum therefore gives
\[
 \sum_{i\in\operatorname{locs}(j,w)}
          \sum_{x\in Y_i}\operatorname{mult}_i(x)\le\dim M_j.
 \tag{LN.21}
\]

The required conversion from $\dim M_j$ to a theta-image rank has
a short complete proof.  Suppose $T:M_j\to H$ and $P:H\to M_j$
are the actual theta maps and, for each nonzero row summand,
$(PT)|_{E_{j,\sigma}}=\lambda_\sigma$ with
$\lambda_\sigma\ne0$.  If $Tv=0$, decompose
$v=\sum v_\sigma$.  Then $0=PTv=\sum\lambda_\sigma v_\sigma$;
independence of the summands forces every $v_\sigma=0$.  Hence
$T$ is injective and $\dim\operatorname{im}T=\dim M_j$.
This is the entire use of \texttt{Conclusion42} here.  It supplies
the actual scalars, not an unproved assertion that dimensions happen
to match.

Thus (LN.21) is \texttt{MultLe}, with right side
$\operatorname{zeroSpaceRankAtPlace}(j,w)$.  The further equality
\[
 \operatorname{zeroSpaceRankAtPlace}(j,w)=d_j,
              \qquad d_j=L.dZ0(j),                     \tag{LN.22}
\]
is the source's \texttt{adm3}.  It follows, for example, by the same
injectivity argument for both this twist and the reference twist,
on the same $M_j$.  Where those scalar identities have not been
supplied, (LN.22) remains an explicit input.  The original
\texttt{GaloisPreconditions.what\textsubscript{2}.rank} supplies it
in the source join.  Neither the existence of a localisation nor
an identification of distinct theta spaces is silently assumed.

\paragraph{LN7: exact weighting, global cut loss, and denominator sum.}
\label{claim:LN7}
Let $\mathcal W$ be the actual finite support of twisted places and
put
\[
 a_w=\frac{e_wf_w\log\ell_w}{[F:\mathbb Q]}\ge0,
 \quad \mathcal E_j=\sum_{w\in\mathcal W}a_w\eta_J(j,w),
 \quad \mathcal V=\sum_{\ell\text{ prime}}\log\ell\,V_\ell.
 \tag{LN.23}
\]
Here $F=D.\mathrm{what}_2.\mathrm{LL}$, $e_w,f_w$ are its actual
ramification index and residue degree, and
$V_\ell=\sum_{1\le e<t}D.\mathrm{what}_2.\mathrm{shaLen}(\ell,e)$.
Thus $a_w$ is literally \texttt{lwP(w)*ramIdx(w)}, and
$\mathcal V$ is the source's \texttt{tsum vvTerm}.  In particular
the ramification factor is present once.

The actual place identification is injective.  For each rational
prime $\ell$, the fundamental degree identity
$\sum_{w\mid\ell}e_wf_w=[F:\mathbb Q]$ gives
\[
              \sum_{w\in\mathcal W,\ w\mid\ell}a_w\le\log\ell.
 \tag{LN.24}
\]
One elementary proof of the degree identity is to compute the
dimension of $\mathcal O_F/\ell\mathcal O_F$ over $\mathbb F_\ell$.
It is $[F:\mathbb Q]$ because $\mathcal O_F$ is a free abelian group
of that rank.  Factorisation of $\ell\mathcal O_F$ and the Chinese
remainder theorem give local quotients of dimensions $e_wf_w$.
This requires no Galois hypothesis on $F$.

The actual tying equation is
$\eta_J(j,w)=\sum_{i\in\operatorname{locs}(j,w)}\eta_i$.
Apply LN5 and LN6, then (LN.24), and finally the prime factorisation
of the positive integer $N_G$.  This proves the exact global loss
\[
 \boxed{\quad\mathcal E_j(D)\le
             \mathcal E_j(\operatorname{cut}D)+d_j\log N_G.\quad}
 \tag{LN.25}
\]
Indeed the remaining scalar sum is bounded by
$\sum_\ell v_\ell(N_G)\log\ell=\log N_G$.  Discarded
localisations contribute zero by LN5.

We also give a direct local-to-global denominator version.  Suppose
there are nonnegative $\nu_\ell,\lambda_\ell$, a positive integer
$h_E$ (the source class number), and, off the two generator primes,
\[
 v\!c_i(x)\le\nu_{\ell_i}+\frac2{\ell_i^2}
                  +\lambda_{\ell_i}+V_{\ell_i}+v_{\ell_i}(h_E)
                       \quad(x\in Y_i).                \tag{LN.26}
\]
At a generator prime assume the actual weighted eta estimate
\[
 \sum_{\substack{w\in\mathcal W\\w\mid\ell}}a_w\eta_J(j,w)
 \le d_j\log\ell\left\{\frac2{\ell^2}+\lambda_\ell
                          +V_\ell+v_\ell(h_E)\right\}.
 \tag{LN.27}
\]
It need not arise from an individual congruence-ideal estimate.
Suppose $\sum_{\ell\in A}\nu_\ell\log\ell\le B_\nu$ for every
finite set of primes $A$, and that the nonnegative series
$\mathcal V$ is summable.  LN5, LN6 and (LN.24) yield
\[
 \boxed{\quad\mathcal E_j\le d_j\left\{
      4+B_\nu+\log h_E+\log N_G+
      \sum_{\ell\in\mathcal P}\log\ell\,\lambda_\ell
                       +\mathcal V\right\}.\quad}     \tag{LN.28}
\]
Here $\mathcal P$ is just the finite set of residue primes of
$\mathcal W$.  Nonnegativity permits extending finite valuation
sums to all prime factors and the finite $V_\ell$ sum to
$\mathcal V$.  The elementary bound
$\sum_{n\ge2}\log n/n^2\le2$ gives the constant $4$; the source
uses the harmless larger $6$.  This preserves coefficient exactly
$d_j$ on $\mathcal V$, an essential feature of the later scalar
join.  No bound on $\mathcal V$ or assertion of its summability is
being proved here.

The source's actual \texttt{ArchimedeanSignatureMatching} in this
branch is the denominator inequality
\[
 v\!c_i(x)\le s_i(x)+\sum_{\lambda\mid\ell_i}b_{i,\lambda}
                 +V_{\ell_i}+v_{\ell_i}(h_E)+B_i+\ell_i^{-2}.
 \tag{LN.29}
\]
Together with $s_i(x)+\ell_i^{-2}\le\nu_{\ell_i}+2\ell_i^{-2}$
and an upper envelope for $\sum b_{i,\lambda}+B_i$, it is exactly
(LN.26).  Its name does not turn (LN.29) into a metric theorem.
The cyclotomic input used by the source is
$\nu_\ell\le v_\ell(QS)+\mathbf1_{\ell\mid QS}v_\ell(tr)$,
with zero at the distinguished excluded primes.  Factorisation then
gives $B_\nu=\log QS+\log(tr)$ directly.  The stated cyclotomic
image inequality, and (LN.29) itself, remain arithmetic inputs.
\paragraph{LN8: the complete numerical prime-envelope estimate.}
\label{claim:LN8}
We now prove the numerical export, including the small-order inert
primes.  This makes the input to LN9 a list of per-prime inequalities,
not a retained global summation theorem.  Use the actual parameter
record of PN16 and its notation from PN1.  Throughout this numerical
argument $n=D.R.n$, $f=D.R.f=2n$, $k=D.R.k$, $m=D.R.m$,
$r=D.R.r=k+m-1$, $t=D.R.t$, $e=D.R.e$ and $o=D.R.o$.
In particular $n$ is the source rank parameter, not a local degree
or the number of selected places.  Set
\[
 \mathcal L=x^\epsilon,\quad b=e^{\mathcal L},\quad
 b'=xb/e^{x^{\epsilon/3}},\quad
 w''=b/(n e^{x^{\epsilon/3}}),\quad S=L.\mathrm{Sc},\quad
 \eta=L.\mathrm{etaSz},\quad d_2=L.\mathrm{de2}.
 \tag{LN.30}
\]
The integer $N_m=k+m(t+1)$ used below is different from $N_G$.
Put
\[
 \begin{split}
 D_*&=nt^2(N_m+t)\mathcal L^2+kn^2(S+t+\mathcal L),\\
 z&=mb/D_*,\qquad \alpha=k\mathcal L/(md_2),\qquad
 H=b+S+\eta+t+2,\qquad \eta_3=\eta+t+2.
 \end{split}                                                   \tag{LN.31}
\]
These are respectively \texttt{Den}, \texttt{lStar}, \texttt{lA},
\texttt{hB}, and \texttt{eta\textsubscript{3}}.  The source's
$w_P$ and $G_P$ obey $w_P+\eta_3\le9b'$ and $G_P\le10b'$.
These follow from the actual window bounds:
$w_B\le3b'+\eta$, $S\le b'$, $t\le\eta\le b'$, with the
constant terms absorbed by $b'\ge1$.

Fix $C\ge1$, independent of $x$ and of every family index, which
bounds the constants and exponents in
\texttt{ParametersPolynomiallyBounded} as well as the local table
constants.  Thus $k,n,t\le x^C$, $d_2\ge x^{-\epsilon}$,
$\log h_E\le Cx^C$, and $\log N_G\le Cx^C\mathcal L$.
The scale has $0<\epsilon<1$.  A sufficiently large threshold,
depending only on the fixed constants and $\epsilon$, gives
\[
 \begin{gathered}
 \mathcal L\ge4,\quad \log x\ge1,\quad
 1\le m\le k,\quad r\le2k,\quad f=2n,\quad n_V\le N_m+t,\\
 \mathcal L\le\alpha<z\le b/\mathcal L,\quad
 (3+\log\alpha)^2\le\mathcal L,\quad
 1\le H\le2b,\quad \mathcal L\le b',\\
 r/\alpha\le2m/\mathcal L,\qquad
 rb'\alpha\log x\le mb,\qquad
 r\eta_3\le4mw'',\qquad
 C(\eta_3+S+\mathcal L)\le w''/8.
 \end{gathered}                                                   \tag{LN.32}
\]
Here are the actual sources of these scalar inequalities, so their
uniformity is not hidden.  Positivity, $m\le k$, $r\le2k$ and
$n_V\le N_m+t$ are rank-record identities.  PN1, PN4--PN6 and PN15
give the scales and the parameter fields
\texttt{rbS}, \texttt{Neta}, \texttt{teta}, \texttt{kS\_le},
\texttt{logeta}, \texttt{logDel}, and \texttt{lstar1}.
The identity $zD_*=mb$, with $D_*\ge kn^2\mathcal L\ge m\mathcal L$,
gives $z\le b/\mathcal L$.  The formulas for $\alpha$ give
$\mathcal L\le\alpha\le x^{C+2\epsilon}$ and
$r/\alpha\le2m/\mathcal L$.
The \texttt{rbS} inequality is
$r(b'+S)\log x\,\alpha\le d_2mb$; discard $S$ and use
$d_2\le1$.  From $N_m\eta\le d_2mw''$, $r\le N_m$ and
$\eta_3\le4\eta$ obtain the penultimate bound.
For the last bound, \texttt{teta} gives $\eta\le w''/x$,
\texttt{kS\_le} and $k\ge\sqrt x$ give $S\le w''/\sqrt x$,
and $n\le x^C$ and exponential domination give
$\mathcal L\le w''/x$.  Thus its left side is at most
$6Cw''/\sqrt x$, which is at most $w''/8$ eventually.
All remaining comparisons in (LN.32) follow from
$\log x=o(x^a)$ for each fixed $a>0$.  For example
$\log\alpha=O(\log x)$, while $D_*$ is subexponential in
$x^{\epsilon/3}$ by the displayed bounds on $\log\eta$ and
$\log\Delta_{\rm cl}$.  Hence $\alpha<mb/D_*$ eventually.
This supplies finite, uniform thresholds, and requires no theorem
about a distribution of primes.

We use the following elementary integer sums, valid for $Y\ge2$
and $A\ge2$ (a finite subset has no larger sum):
\[
 \begin{aligned}
 \sum_{2\le n\le Y}\log n&\le Y\log Y,&
 \sum_{2\le n\le Y}\frac{\log n}{n}&\le(1+\log Y)^2,\\
 \sum_{2\le n\le Y}\frac{\log n}{n-1}&\le2(1+\log Y)^2,&
 \sum_{n\ge A}\frac{\log n}{n^2}&\le\frac{4(\log A+2)}A,\qquad
 \sum_{n\ge2}\frac{\log n}{n^2}\le2.
 \end{aligned}                                                  \tag{LN.33}
\]
For the second bound use $\log n\le\log Y$ and
$\sum_{n\le Y}1/n\le1+\log Y$; the third follows from
$(n-1)^{-1}\le2/n$.  The function $\log u/u^2$ decreases for
$u\ge2$.  Integral comparison with
$\int_A^\infty\log u/u^2\,du=(\log A+1)/A$ proves the fourth,
including a first term and replacing $A$ by its ceiling.  For the
last bound use the first term at $2$ plus the integral from $2$:
$\log2/4+(\log2+1)/2<2$.  The first bound just counts at most
$Y$ nonnegative terms, each at most $\log Y$.

Define
\[
 \begin{split}
 E(\ell)&=C\{1+H/\ell+\log H/\log\ell\},\\
 J(\ell)&=\max\{1,\lceil\log\alpha/\log(\ell-1)\rceil\},\\
 B(\ell)&=H/(\ell-1)^{J(\ell)+1}
                         +(G_P+1)(J(\ell)+1)+\mathcal L,\\
 F(\ell)&=(w_P+\eta_3)/\ell+S+\mathcal L+
                                                \eta_3 I(\ell),\\
 B_0&=Cfnt^2n_V\mathcal L^2.
 \end{split}                                                   \tag{LN.34}
\]
The indicator $I(\ell)$ is $1$ exactly when there is an actual
allowed band parameter $w''\le s\le w_B$ with
\[
 |\ell s-b|\le C(\ell\eta_3+\ell S+b/\ell+s+\ell\mathcal L).
 \tag{LN.35}
\]
This is the source's \texttt{bandI}; replacing a particular
\texttt{band i lambda} by $I(\ell)$ only enlarges an upper bound.
For $\ell\ge2C$, (LN.32) and the upper half of (LN.35) give
\[
 \ell s\le b+\ell s/8+b/2+\ell s/2,
 \qquad w''\le s\le4b/\ell.
 \tag{LN.36}
\]
Consequently
\[
 r\eta_3 I(\ell)\log\ell
 \le16mb\frac{\log\ell}{\ell}
          +\mathbf1_{\ell\le2C}\,4mb\log\ell.       \tag{LN.37}
\]
For the small-prime alternative use $w''\le b'\le b$; if the
indicator is zero there is nothing to prove.

The raw local length inequalities in \texttt{AdmT} imply the
following envelope for an ordinary prime $\ell\ge5$, up to $z$
or in the exceptional set:
\[
 \Lambda(\ell)=B_0+
 \begin{cases}
 (2fm+f)E(\ell),&\ell\text{ split},\\
 4fCrB(\ell),&\ell\text{ inert},\ \ell<\alpha,\ o_\ell>4L_w,\\
 4fCr\{H/\ell^2+F(\ell)\}+2fCmE(\ell),
                    &\ell\text{ inert otherwise}.
 \end{cases}                                                   \tag{LN.38}
\]
Here $o_\ell$ is the least positive exponent with
$\ell^{o_\ell}\equiv\pm1\pmod{2^{e+2}}$.  The class alternatives
are the literal \texttt{SplitK} and \texttt{InertK} away from the
listed special primes.  To check (LN.38) from the individual place
inequalities, use
$\sum_{\lambda\mid\ell}n_\lambda=[E:\mathbb Q]=2f$ and
$\#\{\lambda\mid\ell\}\le2f$.  Thus
$\sum(n_\lambda+1)\le4f$.  The split estimate
$2\sum b_{i,\lambda}\le\sum(2n_\lambda m+1)E(\ell)$ gives
$(2fm+f)E(\ell)$ after division by $2$; the inert estimates give
the two other cases exactly.  This is where the original local
denominator estimates enter, and their truth is not deduced from
the scalar envelope.

We first sum (LN.38) for $5\le\ell\le z$.
The common terms are bounded by
\[
 \sum_{\ell\le z}\log\ell\, B_0\le Cfmb\mathcal L,
 \qquad \sum_{\ell\le z}\log\ell\, E(\ell)
                                      \le11Cb\mathcal L^2.
 \tag{LN.39}
\]
For the first, $znt^2n_V\mathcal L^2\le mb$ and
$\log z\le\mathcal L$.  For the second, expand
$\log\ell E(\ell)=C\{\log\ell+H\log\ell/\ell+\log H\}$.
Use (LN.33), $z\mathcal L\le b$, $H\le2b$ and
$\log H\le\mathcal L+1$.  The three sums are at most
$Cb$, $8Cb\mathcal L^2$ and $2Cb$, respectively.

For the small inert primes of large order, $\ell\ge5$ gives
$(\ell-1)^{J(\ell)}\ge\alpha$ and
$J(\ell)+1\le3+\log\alpha$.  Thus
\[
 B(\ell)\le\frac H{\alpha(\ell-1)}
                      +11b'(3+\log\alpha)+\mathcal L.
 \tag{LN.40}
\]
Multiplying by $r\log\ell$ and summing over $\ell<\alpha$,
the first term is at most
$2(r/\alpha)H(1+\log\alpha)^2\le8mb$.
The other terms are at most
$12rb'\alpha(3+\log\alpha)^2\le12mb\mathcal L$, using
$\mathcal L\le b'$ and $\log x\ge1$.
Multiplication by $4fC$ therefore gives at most
$80Cfmb\mathcal L^2$.
For the medium inert primes, (LN.33) and $r/\alpha\le2m/\mathcal L$
give
\[
 4fCrH\sum_{\ell\ge\alpha}\frac{\log\ell}{\ell^2}
      \le64Cfmb\le64Cfmb\mathcal L^2.                 \tag{LN.41}
\]

The small-order inert primes require a sparsity argument, which
we prove explicitly.  If $\ell$ is odd and
$\ell^o\equiv\pm1\pmod{2^a}$, write $o=2^vu$ with $u$ odd.
Odd powers do not change the valuations of $q-1$ and $q+1$ in
the respective odd-power factors.  If $v>0$, then
$\ell^{2^v}\equiv1\pmod8$, so a congruence to $-1$ modulo
$2^a$ with $a\ge3$ is impossible.  The factorisation into
successive differences of squares gives
\[
 v_2(\ell^{2^v}-1)=v_2(\ell-1)+v_2(\ell+1)+v-1.
 \tag{LN.42}
\]
Indeed the factors $\ell^{2^j}+1$ have valuation $1$ for
$j\ge1$, and exactly one of $\ell-1,\ell+1$ has valuation $1$.
It follows that $\ell\equiv\pm1\pmod{2^{a-v}}$ whenever that
modulus is nontrivial.  The cases $v=0$ or $a\le2$ follow directly
from the odd-power fact and the residues of odd integers.
Taking $a=e+2$, $o=o_\ell\le4L_w$ gives
\[
 \ell\equiv\pm1\pmod M,\qquad
 M=2^{e+2-\lfloor\log_2(4L_w)\rfloor}.               \tag{LN.43}
\]
At the thresholds in use the exponent is nonnegative.

Set $R=n/\mathcal L$, the literal \texttt{LevelStarExponents.R0}.
The field \texttt{radius\_ge\_sqrt\_ratio} gives
$\sqrt{k/m}\le R$, hence $r\le2k\le2mR^2$.
The field \texttt{k\_ge\_mx} gives $R\ge3$ eventually.
The frame field \texttt{ord\_le\_o\_mul\_two\_pow} supplies
$n\le o2^e$, and the size-datum field \texttt{o\_tie} supplies
$1\le o\le x^{\epsilon/100}$ (positivity of $o$ is a frame field).
Finally \texttt{Lw\_le}, together with $P.C_{Lw}\le C$, gives
$L_w\le C\log x$.
Eventually $4o^2\le\mathcal L$,
$2oL_w\le\mathcal L$ and
$(\log R+2)L_w\le\mathcal L$.  The first two and
$R\mathcal L\le o2^e$ imply
$R+1\le2^e/L_w$.  Since
$2^{\lfloor\log_2(4L_w)\rfloor}\le4L_w$, (LN.43) gives
$M\ge2^e/L_w\ge R+1$.  This proves both the claimed sign of
the exponent and the required spacing.  Each of the two progressions
$\pm1\pmod M$, after discarding $1$, starts at $u\ge M-1\ge R$.
For a decreasing $g(t)=\log t/t^2$ on $t\ge3$, interval comparison
gives
\[
 \sum_{h\ge0}g(u+hM)
 \le g(u)+\frac1M\int_u^\infty g(t)\,dt
 \le\frac{\log R}{R^2}+\frac{\log R+2}{RM}.
 \tag{LN.44}
\]
Both functions in the last bound decrease in $R\ge3$.
Combining the two progressions, $r\le2mR^2$ and
$n=R\mathcal L\le oML_w$, we obtain
\[
 r\sum_{o_\ell\le4L_w}\frac{\log\ell}{\ell^2}
 \le4m\log R+
       4mo\frac{(\log R+2)L_w}{\mathcal L}
 \le4m(\log R+o)\le4m\mathcal L^2.                 \tag{LN.45}
\]
Thus this part of (LN.38) costs at most
$32Cfmb\mathcal L^2$.  Counting these primes as an arbitrary
set of primes would lose the essential factor between $r$ and $m$.

For the flat terms, $rb'\le mb$ follows from (LN.32), and
$r(S+\mathcal L)z\le2mb$ follows from the second term of $D_*$.
Use (LN.33) and (LN.37).  Before the factor $4fC$, the four
contributions are at most
\[
 36mb\mathcal L^2,\quad2mb\mathcal L^2,\quad
 64mb\mathcal L^2,\quad4(2C+1)^2mb\mathcal L^2.
 \tag{LN.46}
\]
The last uses
$\sum_{2\le n\le2C}\log n\le(2C+1)^2$.
The extra $E(\ell)$ coefficient in (LN.38) is bounded by
$(3+2C)fm$, since $m\ge1$.  Combining (LN.39)--(LN.46), all
ordinary primes up to $z$ cost at most
\[
 \begin{split}
 &\{C+11C(3+2C)+80C+64C+32C\\
 &\hspace{35mm}+4C[102+4(2C+1)^2]\}\,fmb\mathcal L^2
           \le1000(C+1)^3fmb\mathcal L^2.
 \end{split}                                          \tag{LN.47}
\]
The last inequality follows by expanding the cubic polynomials;
it holds already for every $C\ge0$.

It remains to include the exceptional, small, ramified and generator
primes.  Suppose the actual table supplies the following raw bounds,
exactly the remaining clauses of \texttt{AdmT}:
\[
 \begin{array}{ll}
 \ell>z,\ \ell\notin\mathrm{Exc}:&
              \sum b_{i,\lambda}+B_i=0,\\
 \ell\in\mathrm{small}\text{ or }\ell=3:&
              \sum b_{i,\lambda}+B_i
                   \le Cfmb\log x+Cx^C\mathcal L^2,\\
 \ell=q\text{ or }d:&\sum b_{i,\lambda}+B_i\le2Cfmb,\\
 \ell=g_0\text{ or }g_1:&\lambda_\ell\le Cx^C\eta,\\
 |\mathrm{Exc}|\le(|TT|+2)(S+3\eta+t+8b/z+9),&
                    \ell\in\mathrm{Exc}\Longrightarrow\ell\le3b.
 \end{array}                                                    \tag{LN.48}
\]
The first line is used only off the generator primes and the
special cases in the following lines.  The generator line accompanies
the eta estimate (LN.27), not an assumption about each $v\!c_i(x)$.
In the source \texttt{small} is the set of prime divisors of $QS$,
with $QS\le n^2$ and $2\mid QS$; also $\log q,\log d\le\mathcal L$.
It follows that the small-prime logarithms sum to at most
$2C\log x+\log3$.  The second line therefore costs at most
$4C(C+1)fmb\mathcal L^2$ eventually: use
$\log^2x\le\mathcal L^2$ and absorb
$x^C\log x=o(b)$ in the second term.  The two ramified primes
cost at most $4Cfmb\mathcal L^2$.

For the remaining terms put
$U=x^{C+2}\exp(x^{\epsilon/4})$.  The parameter bounds give
\[
 \begin{gathered}
 S\le4U,\quad \eta,k,n,t,m,\mathcal L\le U,\quad
 f,r\le2U,\quad n_V,N_m+t\le4U^2,\\
 |TT|\le2U,\qquad D_*\le10U^7,\quad b/z\le10U^7.
 \end{gathered}                                                   \tag{LN.49}
\]
The first line follows from \texttt{logeta}, \texttt{logDel},
$S=4k\Delta_{\rm cl}$ and the rank identities; the second follows
by substitution into (LN.31) and the actual \texttt{TT\_card\_le}.
Thus (LN.48) gives $|\mathrm{Exc}|\le388U^8$.
If $\ell>z$, then $\ell>\alpha$, so the large-order small-inert
case of (LN.38) cannot occur.  Since $H\le2b$,
\[
 H/\ell\le20U^7,\quad
 H/\ell^2\le20U^7,\quad
 (w_P+\eta_3)/\ell\le90U^7,\quad
 \log H/\log\ell\le(\mathcal L+1)/\log2.
 \tag{LN.50}
\]
Together with (LN.49), these inequalities give
$\Lambda(\ell)\le K_1(C)U^9$ for an explicit polynomial
$K_1$, for instance $10^4(C+1)^2$.  Also
$\log\ell\le\log(3b)\le3U$ on the exceptional set.
Their whole contribution is at most
$1164K_1(C)U^{18}$.
The two generator primes satisfy
$\log g_k\le C\mathcal L+\log2$ by \texttt{g\_Y} and
$C_{YL}\le C$, so their total cost is at most
$2Cx^C\eta(C\mathcal L+\log2)=O_C(U^3)$.
Finally $4+B_\nu+\log h_E+\log N_G=O_C(U^3)$, using
$B_\nu\le2\log n+\log(tr)$.

For every fixed $a,A>0$, $A U^a/b\to0$: its logarithm is
$\log A+a(C+2)\log x+a x^{\epsilon/4}-x^\epsilon\to-\infty$.
The exceptional, generator and last scalar costs therefore sum to
at most $b\le fmb\mathcal L^2$ at one further uniform threshold.
Combining this with (LN.47)--(LN.48) proves, for \emph{every} finite
set of primes $\mathcal P$,
\[
 4+B_\nu+\log h_E+\log N_G+
        \sum_{\ell\in\mathcal P}\log\ell\,\lambda_\ell
       \le 2000(C+1)^3 fmb\mathcal L^2.               \tag{LN.51}
\]
Here $\lambda$ chooses the generator, small, ramified and ordinary
envelopes in that order, and is zero outside their supported cases.
If cases overlap, bounding by the sum of their nonnegative costs
proves the same estimate.  This completes the source's numerical
export with a harmless different threshold; it preserves the actual
constant $2000(C+1)^3$ and all original local denominator meanings.

\paragraph{LN9: restoration of the original family at one threshold.}
\label{claim:LN9}
Let $\mathrm{Fm}$ be the original \texttt{FamP eps}.  For every
$x$ and $i\in\mathrm{Fm}.I_x$, choose the cut constructed in LN2
when $x$ is beyond the original threshold.  Below that threshold
use the all-points cut; no good-cut assertion is needed there.
The choice is pointwise, but all its subsequent estimates will
have constants independent of $x,i$.  A finite family of indices
is not assumed.

The actual \texttt{gd} construction keeps the family parameter,
index sets and original size constants, and replaces only the
local Hecke points/modules and their congruence losses as follows:
\[
 \begin{array}{c|c}
 \text{original data}&\text{data in the cut family}\\ \hline
 \mathrm{Loc},X_i,T_i,M_i&
 \{i:Y_i\ne\varnothing\},\ S_i,\ T_i|_{S_i},\ M_i/(M_i\cap V_{S_i^c})\\
 \eta_J(j,w)&\displaystyle\sum_{i\in\operatorname{locs}(j,w),\ Y_i\ne\varnothing}
                                                   \eta'_i\\
 F,\mathcal W,e_w,f_w,\mu_{\rm Tw},d_j&\text{the same entries}\\
 f,m,b,\mathcal L,N_c,N_G&\text{the same entries}\\
 \Omega,\mathrm{ldat}_\ell,\mathrm{shaLen},\mathrm{VV},\mathrm{vvTerm}
                                      &\text{the same entries}.
 \end{array}                                                    \tag{LN.52}
\]
The original global theta spaces and maps in the rank calculation
are retained.  This follows by reading the concrete
\texttt{ndCut}, \texttt{layCut}, \texttt{setRaw}, \texttt{ledTr},
and \texttt{extraTr} constructors: \texttt{ndCut} takes the quotient
in (LN.11), while \texttt{setRaw} changes the eta entry and its ties.
The number-field embedding, the \texttt{ldat} and \texttt{sha\_eq}
entries are copied.  The norm choice is the transported original
\texttt{nchTr}, and \texttt{extraTr} transports its accompanying
identities.  Thus the finiteness and summability premises of
\texttt{LinkN} are literally the same propositions in the two
families.  The twisted-place support is also the same, which
proves \texttt{twGood\_gd} with the original threshold.

There is an especially simple bound for the cost of restoration.
Since $\mathrm{GL}_r(\mathbb F_g)$ is a subset of all $r$ by $r$
matrices, $|\mathrm{GL}_r(\mathbb F_g)|\le g^{r^2}$.
Equation (LN.4) gives
\[
       \log N_G\le fr^2(\log g_0+\log g_1).
 \tag{LN.53}
\]
If $k,n\le x^C$, $r\le2k$, $f=2n$ and
$\log g_k\le C\mathcal L+\log2$, then, for $\mathcal L\ge1$,
\[
 \log N_G\le16(C+1)x^{3C}\mathcal L
                     \le C_Nx^{C_N}\mathcal L,
 \quad C_N=\max\{1,3C,16(C+1)\}.
 \tag{LN.54}
\]
This proof needs no class-number bound.  The polynomial rank and
generator-prime logarithm bounds are uniform on the original
family, by its size and parameter records.  Since $\epsilon>0$,
$C_N\log x\le x^\epsilon$ eventually, so $x^{C_N}\le b$.
For every actual lift $1\le j\le m$, $fm\ge1$, and $\mathcal L\ge1$.
Therefore, with the source's exact budget
\[
 \mathcal B=fmb\mathcal L^2+(\log x)^2N_c,
          \qquad N_c\ge0,
 \quad \log N_G\le C_N\mathcal B.                    \tag{LN.55}
\]
Nonnegativity of $N_c$ follows directly from its product/quotient
definition with $x,n,d_2>0$, not from a congruence assertion.

Suppose the cut family has \texttt{LinkN} constant $K$ and
threshold $x_c$, and suppose the multiplicity and rank inputs of
LN6 hold on the original family beyond $x_r$.  At
\[
 x\ge\max\{x_c,x_r,x_{\rm par},x_{\rm exp},
                                      \mathrm{Fm}.x_1+1\},
 \tag{LN.56}
\]
apply the cut estimate with the unchanged \texttt{ShaFinOdd} and
summability premises, then (LN.25) and (LN.55).  For all actual
indices and all allowed $j$,
\[
 \mathcal E_j(D)
 \le d_j\{K\mathcal B+\mathcal V\}+d_j\log N_G
 \le d_j\{(K+C_N)\mathcal B+\mathcal V\}.
 \tag{LN.57}
\]
This proves the complete original-family
\texttt{linkN\_of\_gd} implication.  The coefficient of
$\mathcal V$ stays exactly $d_j$; restoration costs only the
explicit polynomial budget, and never a new cohomological factor.

\paragraph{LN10: complete supplier and precise remaining frontier.}
\label{claim:LN10}
The preceding arguments give the following consolidated supplier.
On one fixed actual family, assume beyond one uniform threshold:
the source-compatible type and spectral data of LN2--LN3; the
row identifications, localisation uniqueness and theta scalar/rank
data of LN6; the parameter record of PN16 and the fixed polynomial
parameter bound in LN8.  Suppose, for every actual $x,i$ satisfying
\texttt{ShaFinOdd}, there is a local denominator table with one
fixed constant $C$ satisfying (LN.29), its cyclotomic scalar bound,
the split/inert length bounds leading to (LN.38), the band condition
(LN.35), and all boundary/exceptional-prime bounds (LN.48).
The table is allowed to depend on $x,i$ and the finiteness premise;
its constant and threshold are not.

If these tables are on the original data, LN1--LN5 supply their
previously independent \texttt{adm1} requirement, and LN6 supplies
\texttt{adm2}.  Apply LN7 and LN8 to the finite set of actual
residue primes.  For every allowed $j$ and summable
\texttt{vvTerm},
\[
 \mathcal E_j\le d_j\{2000(C+1)^3fmb\mathcal L^2+\mathcal V\}
       \le d_j\{2000(C+1)^3\mathcal B+\mathcal V\}.
 \tag{LN.58}
\]
This is \texttt{LinkN} on the original family directly.  In the
source's preferred route the tables are instead on the actual
good-cut data, where the required local regularity is available.
There $Z'_i=Y'_i$, so LN4 alone bounds the cut congruence module
by its cut coordinate ideals; in particular its \texttt{adm1}
holds after adding the nonnegative $v_{\ell_i}(N_G)$ allowance.
The same numerical argument gives \texttt{LinkN} on that cut
family, and LN9 gives (LN.58) on the original family with constant
$2000(C+1)^3+C_N$.

In both cases the order of quantifiers is exactly
\[
 \begin{gathered}
 \exists K,x_0\ \forall x\ge x_0\ \forall i\in\mathrm{Fm}.I_x:\\
 \mathrm{ShaFinOdd}(D_{x,i})\Longrightarrow
 \forall j\in\mathrm{lifts}(D_{x,i}),\\
 \mathrm{Summable}(\mathrm{vvTerm}_{D_{x,i}})
           \Longrightarrow\text{the original-family bound}.
 \end{gathered}
 \tag{LN.59}
\]
Choose $C$ first, increasing it to cover the finitely many original
constant bounds; the inequalities are monotone under this increase
for $x\ge1$.  Then choose the maximum of the finitely many
\emph{uniform} thresholds used above.  Neither step involves a
maximum over the family indices.  All cohomological terms remain
those on the original datum, with precisely the original finiteness
and summability guards.

Here is how this supplier meets the source's local-bound entry
point, without hiding a new analytic assertion in a record name.
A \texttt{Concl} supplies its actual representation family $\rho$,
the ordinary local lengths from its \texttt{HubData}, and a small
prime table $t_f$.  Its combined table has
$s_i(x)=\mathrm{Br.sPi}(\rho_i,x)$.  At primes in
\texttt{small}, set all length entries to zero and put their entire
supplied bound $\sum t_f.\mathrm{len}+t_f.BB$ in $B_i$.
At $3$ outside \texttt{small}, distribute that same total equally
among the finitely many places above $3$; this set is nonempty,
so summing recovers the original total exactly.  At $q,d$ put
the supplied \texttt{midG} in $B_i$.  On the ordinary supported
range use the original Hub lengths and error bound, and put them
zero elsewhere.  Keep the Hub band and exceptional set and set
the generator cost to $Cx^C\eta$ and its scalar error to
$\ell^{-2}$.  This is the literal piecewise \texttt{tab}
construction.  The source's host disjointness and range conditions
make its cases agree with the preceding alternatives.

The pointwise denominator bounds survive this redistribution
because only their sum with $B_i$ occurs in (LN.29).
The split/inert estimates use the untouched ordinary entries;
the small-prime, $q,d$ and generator estimates are precisely their
corresponding supplied \texttt{Concl} clauses.  Large-prime
vanishing and the exceptional-set bounds are unchanged.  Thus a
uniform family of these actual local conclusions, with its fixed
parameter bounds, feeds (LN.59).  This expands the substantive
content of the source's table assembly, local-bound family export
and final \texttt{LinkNFam} restoration.  Their exact declarations
are recorded in the provenance; none is used as an unexplained
numerical theorem.

The remaining leaves are narrower than \texttt{LinkN} and remain
visible.  They are the actual denominator comparison (LN.29),
the cyclotomic scalar inequality, the four local length cases,
the band/exceptional bounds and the boundary-prime estimates,
all on compatible data; and the type/spectral/theta identification
inputs listed in LN2 and LN6.  In the source the first group is
supplied through \texttt{NormChoiceFacts}, local Galois and
integral $p$-adic comparison arguments, including the estimates
at $2$ and $3$.  The label \texttt{Hypotheses\textsubscript{2}}
bundles Fontaine, Artin reciprocity, Bloch--Kato, local duality,
local Euler characteristic, local finiteness and integral
comparison interfaces.  These names do not by themselves certify
that their exact quantified interfaces have already been matched
to classical theorems; any completed LD or other leaf can be
substituted separately in that input ledger.

The source's \texttt{CutGivesAlt} has further requirements.
In addition to the good cut proved in LN2, it demands all norm-choice
facts and \texttt{Tie4} on the transported data.  These are the
literal \texttt{NormChoiceFacts} fields.  Those compatibility assertions
are not manufactured by the quotient construction.  LN1--LN10
close the integral selection, denominator accounting, prime
summation, rank and family-restoration chain from the explicit
local inputs just stated.  Their output is the original
\texttt{LinkN} input in DA2; DA2 still separately requires the
cohomological bound, rank equality and actual metric comparison
specified there.  No claim of a completed RH/GRH proof or of an
unconditional local Galois construction follows from this chapter.
